\documentclass[a4paper,11pt]{article}
\usepackage[T1]{fontenc}
\usepackage{mathpazo}
\usepackage[margin=1.3cm]{geometry}
\usepackage{xcolor,tabularx,titlesec,enumitem,microtype,hyperref}
\definecolor{linkgold}{RGB}{110,84,26}
\hypersetup{colorlinks=true,urlcolor=linkgold,pdftitle={Jinghua Xi -- Curriculum Vitae},pdfauthor={Jinghua Xi}}
\urlstyle{same}\pagestyle{empty}
\setlength{\parindent}{0pt}\setlength{\parskip}{0pt}\setlength{\tabcolsep}{0pt}\raggedright
\titleformat{\section}{\large\bfseries}{}{0pt}{}[\titlerule]
\titlespacing*{\section}{0pt}{7pt}{4pt}
\setlist[itemize]{leftmargin=1.3em,label=\textbullet,topsep=2pt,itemsep=0pt,parsep=0pt,partopsep=0pt}
\setlist[enumerate]{leftmargin=1.8em,label=\textbf{[\arabic*]},topsep=0pt,itemsep=0pt,parsep=0pt}
\newcommand{\cventry}[4]{\par\addvspace{3pt}\begin{tabularx}{\linewidth}{@{}>{\raggedright\arraybackslash}X@{\hspace{1em}}r@{}}\textbf{#1} & #2 \\ \textit{#3} & #4\end{tabularx}\par}
\newcommand{\projectentry}[2]{\par\addvspace{3pt}\begin{tabularx}{\linewidth}{@{}>{\raggedright\arraybackslash}X@{\hspace{1em}}r@{}}\textbf{#1} & #2\end{tabularx}\par}
\begin{document}\fontsize{9.5}{11.2}\selectfont
{\fontsize{28}{32}\selectfont\textbf{Jinghua Xi}}\par\vspace{5pt}
Bonn, Germany \enspace\textbar\enspace \href{mailto:s79jxi@uni-bonn.de}{s79jxi@uni-bonn.de}\par
\href{https://jinghua-xi.github.io/}{jinghua-xi.github.io}
\section{Mathematical Background and Skills}
\textbf{Core areas:} Partial differential equations, functional analysis, and probability theory.\par
\textbf{Programming:} Python; project experience with JavaScript and HTML/CSS.\par
\textbf{Tools used:} Git, SQLite, REST APIs and Ollama.
\section{Education}
\cventry{University of Bonn}{Oct 2025 -- Present}{M.Sc.\ in Mathematics}{Bonn, Germany}
\begin{itemize}
\item \textbf{Current overall GPA: 1.0} (German grading scale; 1.0 is the highest grade).
\item \textbf{Analysis and PDE:} PDE and Functional Analysis (1.0); PDE and Modelling (1.0).
\item \textbf{Probability:} Stochastic Processes (1.3); Optimal Transport Seminar (1.0).
\end{itemize}
\cventry{Southern University of Science and Technology}{Sep 2021 -- Jun 2025}{B.Sc.\ in Mathematics, Fields Honors Program}{Shenzhen, China}
\section{Publications}
\begin{enumerate}
\item C. G. Moreira, J. Xi and Y. Zhang (2025): \textbf{Graphs of continuous but non-affine functions are never self-similar.} \textit{Proc. Amer. Math. Soc.} 153, 2501--2512. \href{https://doi.org/10.1090/proc/17099}{doi:10.1090/proc/17099}.\par\smallskip
\textbf{Main result:} Proved that the graph of a continuous function on $[0,1]$ is self-similar if and only if it is a straight line.
\end{enumerate}
\section{Selected Projects}
\projectentry{AI-Assisted B2B Ordering Prototype}{Nov 2026}
\begin{itemize}
\item Developed and deployed a Python/SQLite ordering prototype integrating self-hosted Qwen3-8B through Ollama on a cloud GPU.
\item Built retrieval over a 100-SKU catalogue, procurement-text parsing, multi-item carts, role-based access and Excel/email workflows.
\item Validated 52 regression tests and implemented one-command SSH deployment with backups and health checks. \href{https://jinghua-xi.github.io/projects/hualian-ordering.html}{Project}
\end{itemize}
\projectentry{LLM Deployment and GPU Inference}{Nov 2026}
Deployed 4-bit quantized Qwen3-8B on an NVIDIA L4 GPU in Google Colab using llama.cpp and CUDA. Used Python to validate an OpenAI-compatible REST API for LLM inference. \href{https://jinghua-xi.github.io/projects/qwen-inference.html}{Project}
\section{Honors and Awards}
\cventry{Rixin Motto Scholarship}{Oct 2024}{Southern University of Science and Technology}{\href{https://newshub.sustech.edu.cn/html/202410/45754.html}{Details}}
\begin{itemize}\item Awarded to 9 students for academic and overall excellence.\end{itemize}
\section{Exchange and Summer Schools}
\cventry{[Exchange] Enhanced Program for Graduate Study}{Feb 2024 -- Jun 2024}{Peking University, China}{\href{https://bicmr.pku.edu.cn/content/show/38-305.html}{Program}}
\begin{itemize}\item Selected as one of 31 participants nationwide.\end{itemize}
\cventry{[Summer School] Geometry}{Aug 2023}{Utrecht University, Netherlands}{\href{https://utrechtgeometrycentre.nl/summer-school/}{Program}}
\cventry{[Summer School] Mathematics of Random Systems}{Aug 2024}{University of Oxford, United Kingdom}{\href{https://www.randomsystems-cdt.ac.uk/oxford-berlin-summer-school-mathematics-random-systems}{Program}}
\cventry{[Summer School] Maximal Functions and Differentiation of Integrals in Euclidean Space}{Jun 2025}{Westlake University, China}{\href{https://its.westlake.edu.cn/info/1017/2920.htm}{Program}}
\section{Leadership Experience}
\cventry{Founder \& Organizer, Undergraduate Mathematics Salon}{Sep 2023 -- Jun 2025}{Southern University of Science and Technology}{\href{https://jinghua-xi.github.io/salon/}{Archive}}
\begin{itemize}
\item Founded a seminar series promoting clear mathematical exposition and discussion among undergraduates.
\item Coordinated monthly events and encouraged students to present and discuss mathematical topics.
\item Featured on the university's official WeChat account. \href{https://mp.weixin.qq.com/s/EpakX8Plu0SW1HP2v0rGkg}{Read more.}
\end{itemize}
\end{document}
